%%%%%%%%%%%%%%%%%%%%%%%%%%%%%%%%%%%%%%%%%%%%%%%%%%%%%%%%%%%%%%%%%%%%%%%%%%%%%%%%%%%%%%%%%%
% Assignment
%%%%%%%%%%%%%%%%%%%%%%%%%%%%%%%%%%%%%%%%%%%%%%%%%%%%%%%%%%%%%%%%%%%%%%%%%%%%%%%%%%%%%%%%%%

\aufgabe
{DE Beispielklausur Aufgabe 1}
{EN Example-exam Assignment 1}

\aufgabenteil{10}
{DE Multiple-Choice Beispiel Aufgabe.}
{EN Multiple-Choice Example Task.}

\fortype{A}{ % ForType kann weggelassen werden für Klausuren mit nur einem Typ
\mcstart[10]{1} % 1 Correct answers, 10 Point each
\mcline{DE Deutscher Text für Choice 1.}{EN English text for choice 1.}{w}
\mcline{DE Deutscher Text für Choice 2.}{EN English text for choice 2.}{f}
\mcline{DE Deutscher Text für Choice 3.}{EN English text for choice 3.}{f}
\mcend
}

\manualText
{DE Manuell eingefügter deutscher Text. Falls man mal außerhalb einer Aufgabe Text in Deutsch und Englisch haben möchte.}
{EN Manually inserted english text. If you want to use text in german or english outside of an assignment.}

\aufgabenteilende

\aufgabenteil{20}
{DE Multiple-Choice Aufgabe mit zwei Auswahlmöglichkeiten pro Zeile.}
{EN Multiple-Choice Task with two Choices per line.}

\fortype{A}{
\mcstarttwo[10]{2} % DE Text, EN Text, Lösung (w/f) für die erste, für die zweite Choice
\mclinetwo
{DE Test1}{EN Test1}{w}
{DE Test2}{EN Test2}{f}
\mclinetwo
{DE Test3}{EN Test3}{w}
{DE Test4}{EN Test4}{f}
\mcendtwo
}

\aufgabenteilende

\aufgabenteil{20}
{DE Beispiel für Eigentschaften Auswahl.}
{EN Example for property selection.}

\fortype{A}{
\propertystart{4}{10} % 4 Kästchen, 1 Punkt pro Zeile. Bei 2 oder 3 Kästchen die weiteren ungenutzten Klammern leer lassen, aber hin schreiben
% Header mit DT Text über der Aufgabe und den Properties (Properties nur EN/DT)
\propertyheader{DE Text}{EN Text}{Prop1}{Prop2}{Prop3}{Prop4}
% Eigentliche aufgabe mit Kästchen.
\propertyline{Beschreibung 1}{Description 1}{w}{f}{f}{w}
\propertyline{Beschreibung 2}{Description 2}{f}{w}{f}{f}
\propertyend
}

\fortype{B}{
\propertystart{4}{1}
\propertyheader{DE Text}{EN Text}{Prop1}{Prop2}{Prop3}{Prop4}
\propertyline{Beschreibung 2}{Description 2}{f}{w}{f}{f}
\propertyline{Beschreibung 1}{Description 1}{w}{f}{f}{w}
\propertyend
}
\aufgabenteilende

\clearpage

\aufgabe{Beispiel Aufgabe 2}{Example Task 2}

\aufgabenteil{20}
{DE Beispiel für eine Freitext-Aufgabe.}
{EN Example for a free-text task.}

% 3 Linien für die Lösung + Musterlösung
\loesung{3}{Loesungstext}

\aufgabenteilende


\aufgabenteil{10}
{DE Inline-Lösungs Beispiel zum Anzeigen von Musterlösungen innerhalb einer Tabelle.}
{EN Inline solution example for showing solutions within tables.}

\begin{center}
ISO/OSI Layer:\\
\begin{tabular}{|c|}
\hline
% Text in der Klausur an der Stelle (z.b. in Tabelle) + Lösung an der Stelle
\lineloesung{L4:~~~~~~~~~~}{Transport} \\
\hline
\lineloesung{L3:~~~~~~~~~~}{Network} \\
\hline
\lineloesung{L2:~~~~~~~~~~}{Data Link} \\
\hline
\lineloesung{L1:~~~~~~~~~~}{Physical} \\
\hline
\end{tabular}
\end{center}
\aufgabenteilende

\aufgabenteil{10}
{DE Ein Bild und die Ausgefüllte Version für die Musterlösung.}
{EN An image to fill out and an image with the correct solution.}

% \bildAufgabe{Caption}{BREITE}{Aufgabenbild Pfad}{Lösungsbild Pfad}
% Caption kann mit \manualText auch in mehreren Sprachen angegeben werden
\bildAufgabe{Image Caption}{0.3\textwidth}{img/bildAfg}{img/bildLsg}
\aufgabenteilende

\aufgabe
{Dolor sit amet}
{Non dolor mundus}

\aufgabenteil{4}
{Freitext}
{Freitext}

\loesung{3}{42 42 42}

\aufgabenteilende


\aufgabenteil{6}
{consectetur adipiscing eli}
{Lorem ipsum dolor sit ame}

\centerline{Beschreibungstect}

\manualText
{Donec iaculis eros ut sodales faucibus?}
{convallis varius ante?}

\loesung{4}{Nunc}

\vspace{0.8cm}

\manualText
{In interdum in lorem non blandit?}
{In hac habitasse platea dic?}

\loesung{4}{4242 \\
euh ueee\\
1000 \\
Wasd}

\aufgabenteilende


\aufgabenteil{1}
{Nunc ultricies aliquet turpis?}
{ Etiam dignissim augue ?}

\mcstarttwo[1]{1}
\mclinetwo
{awgh}{finibus.}{f}
{hrhrjj}{consectetur.}{w}
\mcendtwo
\vspace{-1em}
\aufgabenteilende

\aufgabe
{erat}{semper}


\aufgabenteil{3}
{Donec dapibus eu turpis in egestas}
{ Ut vitae faucibus leo, ut tincidunt lacu}

\propertystart{3}{1}
\propertyheader{aawgh}{reh}{congue}{sapien}{sollicitudin nulla}{}
\propertyline{trjfqe}{consectetur}{w}{f}{w}{}
\propertyline{fwerh}{id}{w}{f}{f}{}
\propertyline{tejrzj}{pharetra}{w}{f}{f}{}
\propertyend

\vspace{-1em}
\aufgabenteilende

\aufgabenteil{3}
{Nullam ut ex aliquam?}
{Vivamus erat ligula?}

\mcstart[1,5]{2}
\mcline{Integer molestie maximus dui, et ornare tortor suscipit sed.}{Vestibulum ante ip.}{f}
\mcline{Maecenas non tempor tortor.}{Aenean congue leo eu sapien pretium aliquam}{f}
\mcline{Proin porta posuere fermentum.}{Morbi facilisis feugiat ex, et placerat arcu semper sodales.}{w}
\mcline{Integer lobortis augue eu.}{Ut consectetur ex id iaculis efficitur. }{f}
\mcline{Aenean a malesuada ante.}{Pellentesque pharetra nec nisl ut fermentum.}{w}
\mcend
\vspace{-1em}

\aufgabenteilende


\aufgabenteil{3}
{Fusce laoreet lacus torto?}
{Proin viverra ligula a pretium semper?}

\mcstarttwo[1,5]{2}
\mclinetwo
{wasd}{gfawg.}{f}
{gaw}{wag.}{w}
\mclinetwo
{hrzh wrg wr}{gawg}{w}
{grgerh erhre}{rhrehr}{f}
\mclinetwo
{wrgh het }{awdasg}{f}
{r zkrz e}{hseghsf}{f}
\mcendtwo
\vspace{-1em}
\aufgabenteilende


\aufgabe{Tabellen}{Tables}

\aufgabenteil{6}
{Table fun\\}
{Table :3\\}

\manualText
{Do table}
{mach mal tabelle}

\begin{center}
\begin{tabular}{|l|l|l|l|l|l|}
\hline
\multicolumn{6}{|l|}{Table} \\ \hline
\scriptsize$F_1$: \lineloesung{~~~~~~}{234} & \scriptsize$F_2$: \lineloesung{~~~~~~}{152} & \scriptsize$F_3$: \lineloesung{~~~~~~}{321} & \scriptsize$F_4$: \lineloesung{~~~~~~}{2412} & \scriptsize$F_5$: \lineloesung{~~~~~~}{215} & \scriptsize$F_6$: \lineloesung{~~~~~~}{23} \\ \hline
\end{tabular}
\end{center}

\aufgabenteilende


\aufgabenteil{5}
{Hier erklärung}
{Ex. Plain.}
\vspace{1em}

\manualText
{Tablle.}
{Just table it}

\begin{center}
\begin{tabular}{|l|l|l|l|l|l|}
\hline
\multicolumn{6}{|l|}{Table} \\ \hline
% Line1
\scriptsize$A_{11}$ \lineloesung{~~~~~}{4242} & \scriptsize$A_{12}$ \lineloesung{~~~~~}{4242} & \scriptsize$A_{13}$ \lineloesung{~~~~~}{4242} & \scriptsize$A_{14}$ \lineloesung{~~~~~}{4242} & \scriptsize$A_{15}$ \lineloesung{~~~~~}{4242} & \scriptsize$A_{16}$ \lineloesung{~~~~~}{4242} \\ \hline
% Line2
\scriptsize$A_{21}$ \lineloesung{~~~~~}{4242} & \scriptsize$A_{22}$ \lineloesung{~~~~~}{4242} & \scriptsize$A_{23}$ \lineloesung{~~~~~}{4242)} & \scriptsize$A_{24}$ \lineloesung{~~~~~}{4242} & \scriptsize$A_{25}$ \lineloesung{~~~~~}{4242)} & \scriptsize$A_{26}$ \lineloesung{~~~~~}{4242} \\ \hline
% Line3
\scriptsize$A_{31}$ \lineloesung{~~~~~}{4242} & \scriptsize$A_{32}$ \lineloesung{~~~~~}{4242} & \scriptsize$A_{33}$ \lineloesung{~~~~~}{4242} & \scriptsize$A_{34}$ \lineloesung{~~~~~}{4242} & \scriptsize$A_{35}$ \lineloesung{~~~~~}{4242} & \scriptsize$A_{36}$ \lineloesung{~~~~~}{4242} \\ \hline
% Line4
\scriptsize$A_{41}$ \lineloesung{~~~~~}{4242} & \scriptsize$A_{42}$ \lineloesung{~~~~~}{4242} & \scriptsize$A_{43}$ \lineloesung{~~~~~}{4242} & \scriptsize$A_{44}$ \lineloesung{~~~~~}{4242} & \scriptsize$A_{45}$ \lineloesung{~~~~~}{4242} & \scriptsize$A_{46}$ \lineloesung{~~~~~}{4242} \\ \hline
\multicolumn{6}{|l|}{Sub Table} \\ \hline
\multicolumn{3}{|l|}{Subtext \manualText{(Deutsch)}{}} & \multicolumn{3}{l|}{Textlol \manualText{(:9)}{}}  \\ \hline
\scriptsize$L_{1}$ 2412 & \scriptsize$L_{2}$ \lineloesung{~~~~~}{4242} & \scriptsize$L_{3}$ \lineloesung{~~~~~}{4242} & \scriptsize$L_{4}$ ~~~2356 & \scriptsize$L_{5}$ ~~~2471 & \scriptsize$L_{6}$ \lineloesung{~~~~~}{4242} \\
\hline
\end{tabular}
\end{center}

\aufgabenteilende

\aufgabe{Platzschinden, groesser 10}{was er geadsgt hat}

\aufgabenteil{3}
{Fusce laoreet lacus torto?}
{Proin viverra ligula a pretium semper?}

\mcstarttwo[1,5]{2}
\mclinetwo
{wasd}{gfawg.}{f}
{gaw}{wag.}{w}
\mclinetwo
{hrzh wrg wr}{gawg}{w}
{grgerh erhre}{rhrehr}{f}
\mclinetwo
{wrgh het }{awdasg}{f}
{r zkrz e}{hseghsf}{f}
\mcendtwo
\vspace{-1em}
\aufgabenteilende


\aufgabenteil{3}
{Fusce laoreet lacus torto?}
{Proin viverra ligula a pretium semper?}

\mcstarttwo[1,5]{2}
\mclinetwo
{wasd}{gfawg.}{f}
{gaw}{wag.}{w}
\mclinetwo
{hrzh wrg wr}{gawg}{w}
{grgerh erhre}{rhrehr}{f}
\mclinetwo
{wrgh het }{awdasg}{f}
{r zkrz e}{hseghsf}{f}
\mcendtwo
\vspace{-1em}
\aufgabenteilende


\aufgabenteil{3}
{Fusce laoreet lacus torto?}
{Proin viverra ligula a pretium semper?}

\mcstarttwo[1,5]{2}
\mclinetwo
{wasd}{gfawg.}{f}
{gaw}{wag.}{w}
\mclinetwo
{hrzh wrg wr}{gawg}{w}
{grgerh erhre}{rhrehr}{f}
\mclinetwo
{wrgh het }{awdasg}{f}
{r zkrz e}{hseghsf}{f}
\mcendtwo
\vspace{-1em}
\aufgabenteilende


\aufgabenteil{3}
{Fusce laoreet lacus torto?}
{Proin viverra ligula a pretium semper?}

\mcstarttwo[1,5]{2}
\mclinetwo
{wasd}{gfawg.}{f}
{gaw}{wag.}{w}
\mclinetwo
{hrzh wrg wr}{gawg}{w}
{grgerh erhre}{rhrehr}{f}
\mclinetwo
{wrgh het }{awdasg}{f}
{r zkrz e}{hseghsf}{f}
\mcendtwo
\vspace{-1em}
\aufgabenteilende

\newpage

\aufgabenteil{3}
{Fusce laoreet lacus torto?}
{Proin viverra ligula a pretium semper?}

\mcstarttwo[1,5]{2}
\mclinetwo
{wasd}{gfawg.}{f}
{gaw}{wag.}{w}
\mclinetwo
{hrzh wrg wr}{gawg}{w}
{grgerh erhre}{rhrehr}{f}
\mclinetwo
{wrgh het }{awdasg}{f}
{r zkrz e}{hseghsf}{f}
\mcendtwo
\vspace{-1em}
\aufgabenteilende


\aufgabenteil{3}
{Fusce laoreet lacus torto?}
{Proin viverra ligula a pretium semper?}

\mcstarttwo[1,5]{2}
\mclinetwo
{wasd}{gfawg.}{f}
{gaw}{wag.}{w}
\mclinetwo
{hrzh wrg wr}{gawg}{w}
{grgerh erhre}{rhrehr}{f}
\mclinetwo
{wrgh het }{awdasg}{f}
{r zkrz e}{hseghsf}{f}
\mcendtwo
\vspace{-1em}
\aufgabenteilende


\aufgabenteil{3}
{Fusce laoreet lacus torto?}
{Proin viverra ligula a pretium semper?}

\mcstarttwo[1,5]{2}
\mclinetwo
{wasd}{gfawg.}{f}
{gaw}{wag.}{w}
\mclinetwo
{hrzh wrg wr}{gawg}{w}
{grgerh erhre}{rhrehr}{f}
\mclinetwo
{wrgh het }{awdasg}{f}
{r zkrz e}{hseghsf}{f}
\mcendtwo
\vspace{-1em}
\aufgabenteilende


\aufgabenteil{3}
{Fusce laoreet lacus torto?}
{Proin viverra ligula a pretium semper?}

\mcstarttwo[1,5]{2}
\mclinetwo
{wasd}{gfawg.}{f}
{gaw}{wag.}{w}
\mclinetwo
{hrzh wrg wr}{gawg}{w}
{grgerh erhre}{rhrehr}{f}
\mclinetwo
{wrgh het }{awdasg}{f}
{r zkrz e}{hseghsf}{f}
\mcendtwo
\vspace{-1em}
\aufgabenteilende

\newpage

\aufgabenteil{3}
{Fusce laoreet lacus torto?}
{Proin viverra ligula a pretium semper?}

\mcstarttwo[1,5]{2}
\mclinetwo
{wasd}{gfawg.}{f}
{gaw}{wag.}{w}
\mclinetwo
{hrzh wrg wr}{gawg}{w}
{grgerh erhre}{rhrehr}{f}
\mclinetwo
{wrgh het }{awdasg}{f}
{r zkrz e}{hseghsf}{f}
\mcendtwo
\vspace{-1em}
\aufgabenteilende


\aufgabenteil{2,5}
{Fusce laoreet lacus torto?}
{Proin viverra ligula a pretium semper?}

\mcstarttwo[2,5]{1}
\mclinetwo
{wasd}{gfawg.}{f}
{gaw}{wag.}{w}
\mclinetwo
{hrzh wrg wr}{gawg}{f}
{grgerh erhre}{rhrehr}{f}
\mclinetwo
{wrgh het }{awdasg}{f}
{r zkrz e}{hseghsf}{f}
\mcendtwo
\vspace{-1em}
\aufgabenteilende


\aufgabenteil{1,5}
{Fusce laoreet lacus torto?}
{Proin viverra ligula a pretium semper?}

\mcstarttwo[1,5]{1}
\mclinetwo
{wasd}{gfawg.}{f}
{gaw}{wag.}{f}
\mclinetwo
{hrzh wrg wr}{gawg}{w}
{grgerh erhre}{rhrehr}{f}
\mclinetwo
{wrgh het }{awdasg}{f}
{r zkrz e}{hseghsf}{f}
\mcendtwo
\vspace{-1em}
\aufgabenteilende


\aufgabenteil{1,5}
{Fusce laoreet lacus torto?}
{Proin viverra ligula a pretium semper?}

\mcstarttwo[1,5]{1}
\mclinetwo
{wasd}{gfawg.}{f}
{gaw}{wag.}{f}
\mclinetwo
{hrzh wrg wr}{gawg}{f}
{grgerh erhre}{rhrehr}{f}
\mclinetwo
{wrgh het }{awdasg}{f}
{r zkrz e}{hseghsf}{w}
\mcendtwo
\vspace{-1em}
\aufgabenteilende

